\input{../tex-inputs/latex-leseansicht-vorspann}

\section[Arthur Schnitzler an Gustav Schwarzkopf, 15. 9. 1899]{L04424 Arthur Schnitzler an Gustav Schwarzkopf, 15. 9. 1899}
\nopagebreak\mylabel{L04424v}
\rehead{ }\normalsize\beginnumbering\briefempfaengerindex{Schwarzkopf, Gustav@\textsc{Schwarzkopf, Gustav}!zzzSchnitzler, Arthur@\emph{von Arthur Schnitzler}!1899-09-152@{15. 9. 1899}|(be}
\toendnotes[C]{\smallbreak\pagebreak[2]}
\correspDesc{Versand  durch Arthur Schnitzler am 15. 9. 1899 in München
\newline{}Übermittlung  am 15. 9. 1899 in München
\newline{}Zustellung  am 18. 9. 1899 in Wien
\newline{}Erhalt  durch Gustav Schwarzkopf am 1[8]. 9. 1899 in Wien}\toendnotes[C]{\smallbreak}
\Standort{DLA, A:Schnitzler, HS.1985.1.1897.}
\physDesc{Bildpostkarte, 336 Zeichen
\newline{}Handschrift: Bleistift, deutsche Kurrent
\newline{}Versand: 1) Stempel: »\nobreak{}\oindex{München@\textbf{München}|pwk}München 1B, 15 Sep 99, 6–7Nm.\nobreak{}«.   2) Stempel: »\nobreak{}\oindex{I., Innere Stadt@\textbf{I., Innere Stadt}|pwk}Wien 1/1 1, 18 9 99, 8–9½V, Bestellt\nobreak{}«. }\toendnotes[C]{\smallbreak}\pstart{}{\pb}Herrn \textsc{Gustav Schwarzkopf}\pend{}\pstart{}Wien\oindex{Wien@\textbf{Wien}|pw}\pend{}\pstart{}\textsc{I. Tiefer Graben 23}\oindex{Wien@\textbf{Wien}!I., Innere Stadt@\textbf{I., Innere Stadt}!Tiefer Graben 23@\textbf{Tiefer Graben 23}|pw}\pend{}{\bigskip}
\pstart
           \centering{}{\pb}\textcolor{gray}{\textbf{München – Lukaskirche\oindex{St. Lukas@\textbf{St. Lukas}|pw}.}}\pend
           \vspace{1em}
\pstart
           \noindent{}{\pb}Herzlichen Gruß! Dank für Ihren \label{K_L04424-1v}\edtext{Brief}{\lemma{\textnormal{\emph{Brief}}}\Cendnote{\textnormal{XXXX Auszeichnungsfehler: Dokument L04292 nicht gefunden.}}}\label{K_L04424-1}. Die
               Radtour fällt natürlich aus (Hochwaſſer \textsc{etc.}) aber das
               übrige iſt in Ordnung.– Wahrſcheinlich fahr ich dann \label{K_L04424-2v}\edtext{nach Frankfurt\oindex{Frankfurt am Main@\textbf{Frankfurt am Main}|pw}}{\lemma{\textnormal{\emph{nach Frankfurt}}}\Cendnote{\textnormal{Vgl. A. S.: \emph{Wiener Schnitzler}, 19. 9. 1899.}}}\label{K_L04424-2} zu Goldmann\pwindex{Goldmann, Paul 31.\,1.\,1865 Breslau – 25.\,9.\,1935 Wien@\textsc{Goldmann, Paul} (31.\,1.\,1865 Breslau – 25.\,9.\,1935 Wien)|pw}.
                  \label{K_L04424-3v}\edtext{Heute{ }ſeh ich mir die
                  Vermächtnis\pwindex{Schnitzler, Arthur 15. 5. 1862 Wien – 21. 10. 1931 ebd.@\textsc{Schnitzler, Arthur} (15. 5. 1862 Wien – 21. 10. 1931 ebd.)!Vermächtnis. Schauspiel in drei Akten@\strich\emph{Das Vermächtnis. Schauspiel in drei Akten}|pw}première\eventindex{Münchner Schauspielhaus@\textbf{Münchner Schauspielhaus}!Premiere von Das Vermächtnis, 15.9.1899@Premiere von Das Vermächtnis, 15.9.1899|pw}}{\lemma{\textnormal{\emph{Heute … Vermächtnispremière}}}\Cendnote{\textnormal{Premiere von Das Vermächtnis, 15.9.1899. In: A. S.: \emph{Kulturveranstaltungen}, \url{https://schnitzler-kultur.acdh.oeaw.ac.at/pmb209099.html}.}}}\label{K_L04424-3}{ }\textsc{incognito} von der Gallerie\oindex{Münchner Schauspielhaus@\textbf{Münchner Schauspielhaus}|pwv} an. – \label{K_L04424-4v}\edtext{Geſtern die  3 Einakter\pwindex{Schnitzler, Arthur 15. 5. 1862 Wien – 21. 10. 1931 ebd.@\textsc{Schnitzler, Arthur} (15. 5. 1862 Wien – 21. 10. 1931 ebd.)!grüne Kakadu – Paracelsus – Die Gefährtin. Drei Einakter@\strich\emph{Der grüne Kakadu – Paracelsus – Die Gefährtin. Drei Einakter}|pwv}\eventindex{Residenztheater München@\textbf{Residenztheater München}!Aufführung von Paracelsus, Die Gefährtin, Der grüne Kakadu, 14.9.1899@Aufführung von Paracelsus, Die Gefährtin, Der grüne Kakadu, 14.9.1899|pwv}}{\lemma{\textnormal{\emph{Gestern die  3 Einakter}}}\Cendnote{\textnormal{Siehe A. S.: \emph{Kulturveranstaltungen}, \url{https://schnitzler-kultur.acdh.oeaw.ac.at/pmb32562.html}.}}}\label{K_L04424-4}, amuſirte mich ſehr  gut.\pend
           \pstart \spacefill\mbox{A.}\pend{}\selectlanguage{ngerman}\endnumbering\briefempfaengerindex{Schwarzkopf, Gustav@\textsc{Schwarzkopf, Gustav}!zzzSchnitzler, Arthur@\emph{von Arthur Schnitzler}!1899-09-152@{15. 9. 1899}|)be}\mylabel{L04424h}
\begin{anhang}
\end{anhang}\newcommand{\dateiname}{L04424}\newcommand{\titel}{Arthur Schnitzler an Gustav Schwarzkopf, 15. 9. 1899}\newcommand{\editorInnen}{Herausgegeben von Jahnke, SelmaMüller, Martin Anton}\input{../tex-inputs/latex-leseansicht-abspann}
